%!TEX TS-program = pdflatex
%!TEX encoding = UTF-8 Unicode

\usepackage{etoolbox}
\AtBeginEnvironment{Shaded}{\scriptsize}

\makeatletter
% Save the original include graphics command
\let\old@includegraphics\includegraphics

% Redefine it to check if a width was provided
\renewcommand{\includegraphics}[2][]{%
  \centering
  \ifstrempty{#1}{%
    % No optional arguments (like width) provided: set default to 0.5\textwidth
    \old@includegraphics[width=0.5\textwidth]{#2}%
  }{%
    % Optional arguments provided: use them (this preserves your {width=100%})
    \old@includegraphics[#1]{#2}%
  }%
}
\makeatother

% TUM Styleguide: https://portal.mytum.de/corporatedesign/regeln/index_styleguide
\definecolor{normalblue}{RGB}{0, 101, 189}
\definecolor{darkblue}{RGB}{0, 51, 89}
\definecolor{lighterblue}{RGB}{34, 111, 167}
\definecolor{darkerblue}{RGB}{0, 37, 73}
\definecolor{ivory}{RGB}{218, 215, 203}
\definecolor{matbluegreylight}{RGB}{236, 239, 241}

\setbeamercolor{structure}{fg=normalblue} % Subtitles and bulletpoints
\setbeamercolor{frametitle}{fg=normalblue} % Slide title
\setbeamercolor{title}{fg=normalblue} % Presentation title
\setbeamercolor{palette tertiary}{bg=normalblue, fg=white} % Left decoration color
\setbeamercolor{palette secondary}{bg=darkblue, fg=white} % Mid bottom decoration color
\setbeamercolor{palette primary}{bg=darkerblue, fg=white} % Right decoration color

\setbeamercolor{author in head/foot}{fg=white, bg=normalblue} % Left deco
\setbeamercolor{title in head/foot}{fg=white, bg=darkblue} % Mid deco
\setbeamercolor{date in head/foot}{fg=white, bg=darkerblue} % Right deco

\setbeamertemplate{blocks}[rounded][shadow=false]
\setbeamercolor{block title}{fg=black,bg=matbluegreylight}
\setbeamercolor{block body}{fg=black,bg=matbluegreylight}
\setbeamercolor{block title}{use=structure,fg=structure.fg,bg=matbluegreylight}
\setbeamercolor{block body}{parent=normal text,use=block title,bg=matbluegreylight}

\usepackage{mathrsfs}
\usepackage{amsmath}
